\documentclass[10pt,a4paper]{amsart}
\usepackage[margin=19mm]{geometry}
\usepackage{amsmath,amssymb,mathtools}
\usepackage[scr=rsfs,cal=euler]{mathalfa}
\usepackage{xfrac}
\usepackage[unicode,hidelinks,bookmarks=false]{hyperref}
\newcommand{\ie}{i.e.\ }
\newcommand{\cf}{cf.\ }
\newcommand{\resp}{respectively\ }
\newcommand{\Subsection}{\subsection}
\providecommand{\nobreakdash}{\nobreak\hbox{-}}
\setlength{\emergencystretch}{2em}
\multlinegap0pt
\usepackage{amssymb}
\usepackage{mathtools}
\usepackage{dsfont}
\usepackage[shortlabels]{enumitem}
\newtheorem{lemma}{Lemma}
\newcommand{\C}{\mathbb C}
\newcommand{\be} {\begin{equation}}
\newcommand{\bea} {\begin{eqnarray}}
\newcommand{\beas}{\begin{eqnarray*}}
\newcommand{\cont}{\mathcal C}
\newcommand{\ee} {\end{equation}}
\newcommand{\eea} {\end{eqnarray}}
\newcommand{\eeas}{\end{eqnarray*}}
\newcommand{\ltt}{\left(}
\newcommand{\rtt}{\right)}
\title{Dual realizations of Bergman spaces on strongly convex domains}
\author{Agniva Chatterjee}
\date{Source: arXiv:2506.02913v2 (CC BY 4.0)}
\begin{document}
\maketitle

\noindent\emph{Source and licence.} Dual realizations of Bergman spaces on strongly convex domains. Version arXiv:2506.02913v2 (CC BY 4.0), distributed by arXiv under the Creative Commons Attribution 4.0 International licence, http://creativecommons.org/licenses/by/4.0/, as recorded in the arXiv licence metadata for this exact version and shown on the abstract page https://arxiv.org/abs/2506.02913v2. Full licence notice: \texttt{licenses/arxiv-2506.02913-CC-BY-4.0.txt}.\par\smallskip

\noindent\textit{Editorial scope.} Counted unit: the article's lemma Minkwski (source0.tex lines 383--386). The article's complete original proof of it is delivered with it (source0.tex lines 387--412, one \texttt{proof} environment of 26 lines). No mathematics is written, repaired or re-derived: the statement and the proof are the article's own lines, in the article's own notation, and no retained line is substituted or re-flowed.  No further source-local context is needed: every cross-reference of every retained line resolves inside the two delivered spans.  Nothing was omitted from any retained span, so the package carries no omission record.  No retained line contains a citation, so the wrapper carries no bibliography and no bibliography entry.  No retained line contains a figure environment, an image-inclusion command or drawing code, so no image is delivered and the package declares no asset dependency.  Editorial formatting only, in addition: an AMS \texttt{amsart} preamble that restates the article's own declarations and macros used by the retained lines, namely the environments \texttt{lemma} on separate counters, together with the standard packages those lines need.  The retained environments print in this document as Lemma 1 in order of appearance, whereas the article prints the same items with the numbering of its own full document.  The complete native source of the article as distributed in the arXiv e-print for this exact version is bundled unchanged as \texttt{sources/179235/source0.tex}.
\begin{lemma}\label{le:str cvx of Minkwski}
Let $D\subset\C^n$ be a bounded $\cont^2$-smooth strongly convex domain containing the origin. Then $\rho_D(\zeta)=m_D^2(\zeta)-1$ is a strongly convex function on $\C^n\setminus\{0\}$, i.e.,
$\operatorname{Hess}(\rho_D)(\zeta)$ is positive definite, for all $\zeta\in\C^n\setminus\{0\}$.
\end{lemma}

\begin{proof}
Let $\ltt \zeta,w\rtt\in \ltt\C^n\setminus\{0\}\rtt^2$. Writing $w$ in real co-ordinates as $w=\ltt u_1,u_2,\cdots,u_{2n}\rtt$, we get that,
\beas
w^T \cdot \operatorname{Hess}(\rho_D)(\zeta)\cdot w &=& 
\sum_{j,k=1}^{2n} \frac{\partial^2\rho_D}{\partial x_j \partial x_k}(\zeta) u_ju_k\\
&=&2\sum_{j,k=1}^{2n} \ltt\frac{\partial m_D}{\partial x_j}(\zeta)\frac{\partial m_D}{\partial x_k}(\zeta)+m_D(\zeta) \frac{\partial^2 m_D}{\partial x_j \partial x_k}(\zeta)\rtt  u_ju_k \\
&=&2\ltt \sum_{j=1}^{2n}\frac{\partial m_D}{\partial x_j}(\zeta) u_j\rtt^2+2m_D(\zeta) \sum_{j,k=1}^{2n}\frac{\partial^2 m_D}{\partial x_j \partial x_k}(\zeta) u_ju_k.
\eeas
Since $m_D$ is positively $1$-homogeneous, the right-hand side is $0$-homogeneous in $\zeta$,  i.e., $f(t\zeta)=f(\zeta),$ for all $\zeta\in\C^n\setminus\{0\}$ and $t>0$. 
Thus, it suffices to show that
\be\label{eq:claim str cvx}
2\ltt \sum_{j=1}^{2n}\frac{\partial m_D}{\partial x_j}(\zeta) u_j\rtt^2+2m_D(\zeta) \sum_{j,k=1}^{2n}\frac{\partial^2 m_D}{\partial x_j \partial x_k}(\zeta) u_ju_k>0,\quad\forall \zeta\in bD\,\text{and}\,\forall w\in\C^n\setminus\{0\}.
\ee
Now, as $D$ is convex, $m_D$ is a convex function on $\C^n$, which implies that 
\be\label{eq:mink cvx}
\sum_{j,k=1}^{2n}\frac{\partial^2 m_D}{\partial x_j \partial x_k}(\zeta) u_ju_k\geq 0,\quad\text{for all } \zeta\in bD\,\text{and } w\in\C^n\setminus\{0\}.
\ee
Furthermore, as $\rho_D$ is a defining function of $D$, it follows that for any $\zeta\in bD$ and $w\in \C^n\setminus\{0\}$, 
\bea\label{eq:tangent}
\text{if }
\sum_{j=1}^{2n}\frac{\partial m_D}{\partial x_j}(\zeta) u_j=0,
\quad \text{then }
\sum_{j,k=1}^{2n}\frac{\partial^2 m_D}{\partial x_j \partial x_k}(\zeta)u_ju_k>0.
\eea
Thus, \eqref{eq:claim str cvx} follows  from \eqref{eq:mink cvx} and \eqref{eq:tangent}.
\end{proof}

\end{document}
